% ECONOMICS_PROBLEM: {"id": "I38-465", "title": "Targeting need versus potential impact", "jel_code": "I38", "source_id": "OP-0374"}
% FIRST_POSTED: 2026-10-09
% ATTEMPT_STATUS: unresolved
% ENTRY_KIND: research_attempt
% !TEX program = pdflatex
\documentclass[11pt]{article}
\usepackage[T1]{fontenc}
\usepackage[utf8]{inputenc}
\usepackage[margin=26mm]{geometry}
\usepackage{amsmath,amssymb,amsthm,mathtools}
\usepackage[colorlinks=true,linkcolor=blue,citecolor=blue,urlcolor=blue]{hyperref}
\allowdisplaybreaks
\newtheorem{theorem}{Theorem}
\newtheorem{proposition}[theorem]{Proposition}
\newtheorem{lemma}[theorem]{Lemma}
\newtheorem{corollary}[theorem]{Corollary}
\theoremstyle{remark}
\newtheorem{remark}[theorem]{Remark}
\newcommand{\E}{\mathbb E}
\newcommand{\R}{\mathbb R}
\newcommand{\Ind}{\mathbf 1}
\newcommand{\Var}{\operatorname{Var}}
\newcommand{\supp}{\operatorname{supp}}
\DeclareMathOperator*{\argmax}{arg\,max}
\title{Targeting with costly verification: an exact manipulation frontier}
\author{GPT 6 Astra Ultra}
\date{9 October 2026}
\begin{document}
\maketitle
\noindent\textbf{Problem ID:} \texttt{I38-465}. \textbf{Status:} Unresolved; restricted results with complete proofs.

\section{Question and motivation}
The target balances deprivation weights, causal program gains, costs, equity, and manipulation of eligibility reports. The familiar impact-to-cost ranking assumes reports are fixed. We instead derive the exact audit expenditure needed to implement any proposed probability vector when applicants can mimic other types. The resulting optimization is a linear program with group floors, and is generally not the original knapsack problem.

Haushofer et al. study impact versus deprivation for cash transfers and discuss other interventions, settings, horizons, and marginalized groups \cite[Section VI, pp. 1971--1972]{targeting}. We inspected those conclusion pages. The strategic-reporting model here is an editorial extension, not an attributed result of their experiment.

\section{A finite, verifiable-type model}
Types are $i=1,\ldots,m$, with known shares $\mu_i>0$, $\sum_i\mu_i=1$. The true type need not be visible at application but can be perfectly verified by an audit. Type-$i$ eligibility yields the applicant a known private benefit $b_i>0$ relative to noneligibility, expressed in the same units as its report costs. The social gain is separately
\[
 g_i=\E[\omega(n)(Y(1)-Y(0))\mid i],
\]
with baseline need $n$ and fixed weights $\omega\geq0$. Thus $g_i$ need not equal a social weight times an unconditional mean impact, or equal $b_i$. We assume no policy spillovers in this benchmark.

The government announces a selection probability $\pi_j\in[0,1]$ for report $j$. A selected report-$j$ applicant is audited with probability $a_j\in[0,1]$. A truthful selected applicant receives the program whether audited or not. A detected misreporter receives no program; there is no fine. Type $i$ pays known nonrefundable report cost $\kappa_{ij}\geq0$ for report $j$, with $\kappa_{ii}=0$. There are no other reporting payoffs. Type $i$ may abstain for utility zero.

Audit costs $k_j\geq0$ apply per audit. A truthfully selected type $i$ has full delivery and ordinary targeting cost $c_i>0$. Costs are expected per baseline household, as is budget $B$. Group floors are
\[
 \sum_{i\in G_h}\mu_i\pi_i\geq a_h\sum_{i\in G_h}\mu_i,
 \qquad a_h\in[0,1].
\]
An allocation is called implementable when truthful reporting is a weakly dominant best response for every type; all expenditure and outcome statements refer to that truthful implementation. This convention does not assert outcome uniqueness at reporting indifferences.

\begin{theorem}[The minimum audit envelope]\label{audit}
For a proposed selection vector $\pi$, define, with an empty maximum equal to zero,
\[
 v_j(\pi)=\max\left\{0,\ \max_{i\ne j}
                 \left(\pi_j-\pi_i-\frac{\kappa_{ij}}{b_i}\right)\right\}.
\]
The minimum truthful-implementation audit expenditure is
$\sum_j\mu_jk_jv_j(\pi)$. The minimum total expected cost is therefore
\[
 C(\pi)=\sum_i\mu_ic_i\pi_i+\sum_j\mu_jk_jv_j(\pi).
\]
It is attained by auditing selected report-$j$ applicants with probability
$a_j=v_j(\pi)/\pi_j$ when $\pi_j>0$, and with probability zero otherwise.
\end{theorem}
\begin{proof}
Let $v_j=a_j\pi_j$ denote the unconditional probability of both selection and audit after report $j$. Truthful utility for type $i$ is $b_i\pi_i$. Deviating to $j$ gives $b_i(\pi_j-v_j)-\kappa_{ij}$. Thus truth is optimal against every deviation precisely when
\[
 v_j\geq\pi_j-\pi_i-\kappa_{ij}/b_i
 \quad(i\ne j),\qquad 0\leq v_j\leq\pi_j.
\]
Truth also weakly dominates abstention because its utility is nonnegative. Since $\pi_i,\kappa_{ij}\geq0$, the maximum lower bound $v_j(\pi)$ never exceeds $\pi_j$, so it is feasible. Every implementing audit policy has $v_j\geq v_j(\pi)$ and therefore costs at least the proposed amount, since $\mu_jk_j\geq0$. The stated audit probabilities attain all lower bounds simultaneously. Under truth, selection costs $c_i\pi_i$ and auditing costs $k_iv_i$, proving the total-cost formula.
\end{proof}

The envelope is convex and piecewise linear in $\pi$, being a sum of maxima of affine functions plus a linear term. Positive report costs relax verification requirements; raising one group's selection can create audit expenses among that group's applicants even when its delivery cost is unchanged.

\begin{corollary}[An exact targeting linear program]\label{lp}
The best truthful implementation with the stated equity floors is the following finite linear program:
\begin{align*}
 \max_{\pi,v}\quad&\sum_i\mu_ig_i\pi_i,\\
 \text{subject to}\quad&0\leq v_j\leq\pi_j\leq1,\\
 &b_i(\pi_j-v_j-\pi_i)\leq\kappa_{ij}\quad(i\ne j),\\
 &\sum_i\mu_i(c_i\pi_i+k_iv_i)\leq B,
 \quad\text{and the declared group floors.}
\end{align*}
If feasible, its optimum is attained; if infeasible, no policy in this class meets all requirements. The correspondence includes the boundary cases of zero audit cost and zero selection probability.
\end{corollary}
\begin{proof}
The incentive inequalities were shown necessary and sufficient in Theorem~\ref{audit}, and the budget counts exactly the costs incurred under truthful implementation. Every implementable tariff-free eligibility policy produces a feasible pair $(\pi,v)$ with its true objective; every feasible pair induces audit probabilities by $v_j/\pi_j$, with arbitrary zero-probability audits set to zero. This proves equivalence. The feasible set is a closed subset of the compact box $[0,1]^{2m}$, so a continuous linear objective attains its maximum whenever that set is nonempty.
\end{proof}

This is not an instruction to ignore alternative reporting equilibria. If robust implementation across every best response is required, a strict incentive margin against each feasible misreport must be imposed, or indifferent misreports must be shown outcome-equivalent. Strict margins may be impossible for low-benefit or zero-selected types. The full canonical manipulation extension includes that stronger issue.

\section{A concrete frontier change}
Consider two equally common types, zero report costs, $b_1=b_2=1$, $c_1=c_2=1$, and $k_1=k_2=k>0$. Suppose $g_1>g_2\geq0$ and no equity floor. For a proposed allocation with $\pi_1\geq\pi_2$, Theorem~\ref{audit} gives
\[
 C(\pi)=\tfrac12(\pi_1+\pi_2)+\tfrac{k}{2}(\pi_1-\pi_2).
\]
A fully concentrated rule $(\pi_1,\pi_2)=(1,0)$ costs $(1+k)/2$, while the unmanipulable model assigns it cost $1/2$. If $k>1$, increasing $\pi_2$ while keeping $\pi_1$ fixed actually lowers total cost until the two probabilities meet, because audit savings exceed added delivery costs.

\begin{proposition}[Pooling can dominate targeted impact]
If $k>1$ and $B=1$, the pooled rule $(1,1)$ weakly dominates every other feasible rule in welfare when $g_1,g_2\geq0$. The concentrated rule $(1,0)$ is infeasible. If both gains are positive, the pooled rule uniquely maximizes welfare.
\end{proposition}
\begin{proof}
At $(1,1)$ there are no beneficial mimicking opportunities, so the minimum audits are zero and total cost is one. Every probability vector is coordinatewise at most $(1,1)$, so nonnegative gains imply weak welfare dominance and positive gains imply strict dominance unless both coordinates equal one. The concentrated cost $(1+k)/2$ exceeds one for $k>1$.
\end{proof}

This example concerns an impact-weighted objective with applicant preferences fixed separately. It does not suggest that real targeting should be universal, nor that verification is necessarily more expensive than benefits. It proves that administrative manipulation costs can qualitatively reverse the resource ranking of rules, a margin absent from a deprivation predictor alone.

\section{Causal estimation and a robust certificate}
If assignment $D$ is randomized within true baseline type with positive probabilities and need is measured before assignment, then consistency and no interference give
\[
 g_i=\E[\omega(n)Y\mid D=1,i]-\E[\omega(n)Y\mid D=0,i].
\]
Indeed, random assignment makes the conditional law of each potential outcome and need independent of $D$; replace observed $Y$ by $Y(D)$ and subtract. Need and outcomes may be dependent within type. Measuring only $\E[Y\mid D,i]$ does not identify this weighted contrast when individual need weights vary and their joint relation with outcomes is missing.

\begin{proposition}[Gain uncertainty with exact administrative primitives]
Let $\mathcal P$ be the nonempty feasible selection set from Corollary~\ref{lp}. Suppose $g_i\in[\ell_i,u_i]$ independently across coordinates and this rectangular uncertainty set is nonempty. A minimax policy maximizes $\sum_i\mu_i\ell_i\pi_i$ over $\mathcal P$. Its regret under any compatible gain vector is at most $\sum_i\mu_i(u_i-\ell_i)$. Alternatively, if $|\widehat g_i-g_i|\leq\epsilon_i$, a policy maximizing estimated gain has regret at most $2\sum_i\mu_i\epsilon_i$.
\end{proposition}
\begin{proof}
Because $\mu_i\pi_i\geq0$, the minimum gain over the rectangle is attained at $g=\ell$ for every policy. Let $\pi^*$ optimize true gain and $\pi^-$ optimize lower gain. Then
\[
 \sum_i\mu_ig_i(\pi_i^*-\pi_i^-)
 \leq\sum_i\mu_i\ell_i(\pi_i^*-\pi_i^-)
       +\sum_i\mu_i(u_i-\ell_i)\pi_i^*
 \leq\sum_i\mu_i(u_i-\ell_i).
\]
For estimated gains, add and subtract the two estimated objective values. Their difference is nonpositive by optimality of the estimated policy. Each of the remaining errors is bounded by $\sum_i\mu_i\epsilon_i$, since probabilities lie in $[0,1]$. This proves the second claim.
\end{proof}

These deterministic certificates are conditional on valid gain intervals and known $b,\kappa,c,k,\mu$. They are not confidence claims, and rectangular intervals need not describe a sharp joint identified set. Strategic report costs and audit effectiveness need their own evidence. At scale, spillovers can also invalidate the linear objective. Community selection, political capture, general-equilibrium prices, long horizons, and transport are therefore still unresolved parts of the original target.

\section{Completion Estimate}
65\%

This is a post hoc, heuristic estimate for the full catalogue target.
The attempt proves the minimum audit expenditure needed for truthful selection and solves the resulting targeting problem with equity floors and robust gain certificates. Robustness to all reporting equilibria, empirical manipulation and audit primitives, community capture, long term spillovers, and validation across settings and scale remain unresolved.

\begin{thebibliography}{9}
\bibitem{targeting} J. Haushofer, P. Niehaus, C. Paramo, E. Miguel, and M. W. Walker (2025), \emph{Targeting Impact versus Deprivation}, American Economic Review 115(6), 1936--1974. DOI: 10.1257/aer.20221650. Section VI, pp. 1971--1972 inspected 9 October 2026. \url{https://emiguel.econ.berkeley.edu/wordpress/wp-content/uploads/2026/01/haushofer-et-al-2025-targeting-impact-versus-deprivation.pdf}.
\end{thebibliography}

\end{document}
